% titel: Blatt 6, A6: Resonanzformeln von Stickstoffverbindungen
% skript: Mortimer Kap. 8 (chemie: Atombindung, Resonanz)

Resonanz beschreibt Moleküle, deren Elektronenlokalisation nicht durch eine
einzige Lewis-Struktur darstellbar ist: mehrere Strukturen mit gleicher
Atomposition, aber verteilten Doppelbindungen/freien Paaren werden von der
realen Elektronenwolke gemittelt (Mesomerie).

\begin{werkzeugkasten}[Resonanzformeln schreiben]
\begin{itemize}
  \item Nur die Position der $\pi$-Elektronen (Doppelbindungen) und freier
        Paare \textit{shuften sich} --- Atome bleiben ortsfest.
  \item Jede Resonanzstruktur muss Lewis-gültig sein (Oktetts möglichst
        erfüllt, Formalladungen notiert).
  \item Die tatsächliche Struktur ist die Überlagerung (Resonanzhybrid);
        energetisch niedrige Strukturen (Formalladung nahe 0, volle Oktetts)
        tragen am meisten bei.
  \item Stickstoff trifft auf N mit Resonanz vor allem in $\mathrm{NO}$,
        $\mathrm{NO_2}$, $\mathrm{NO_3^-}$, $\mathrm{N_2O}$.
\end{itemize}
\textit{Skript:} Mortimer Kap. 8 (Resonanz, Mesomerie).
\end{werkzeugkasten}

\begin{achtung}[Darauf achten]
\begin{itemize}
  \item Formalladungen sorgfältig setzen (z.\ B. $\mathrm{N^+}$ und $\mathrm{O^-}$
        in $\mathrm{NO_2^-}$) --- die Summe muss die Ionenladung treffen.
  \item $\mathrm{NO_3^-}$ hat drei \textit{äquivalente} Resonanzstrukturen (eine
        $\mathrm{N=O}$-Doppelbindung je O), nicht drei verschiedene Moleküle.
  \item Elektronen zählen, bevor man Doppelbindungen verteilt --- so entstehen
        keine falschen Strukturen wie $\mathrm{N}$ mit 5 Bindungen.
\end{itemize}
\end{achtung}
